\documentclass[11pt]{article}
\usepackage[a4paper,margin=26mm]{geometry}
\usepackage{amsmath,amssymb,amsthm,hyperref}
\newtheorem{theorem}{Theorem}
\newtheorem{lemma}{Lemma}
\newcommand{\E}{\mathbb E}
\newcommand{\Pol}{\operatorname{Pol}}
\title{A quadratic lower bound for every permutation-fair die}
\author{Research proof: three independent full mathematical reviews completed}
\date{30 September 2026}
\begin{document}\maketitle
\noindent\textbf{Status.} Three independent full mathematical reviews
completed, with no gap identified; this is not formal verification.
The proof combines two independently reviewed inputs with an
all-orders Gaussian estimate.

\begin{theorem}\label{main}
There is an absolute $c>0$ such that every individual die in a tie-free
permutation-fair collection of $n$ equiprobable finite dice has at least
$cn^2$ faces. More generally, for any monotone mixed-moment tensor with
$m$ columns, its first and last row weights are at most $C/m^2$.
Consequently the smallest possible size of one distinguished die,
with all other dice unrestricted, is $\Theta(n^2)$.
\end{theorem}
The last assertion uses the independently reviewed $O(n^2)$ construction
in \path{polynomial_exceptional_die.tex}. No claim about literature
novelty or about the optimal total number of faces is made here.

\section*{Notation and reviewed inputs}
Assume $p_j(q)\in[0,1]$ are nondecreasing in $q$, the positive row weights
$w_q$ sum to one, and
\[
 \E_w\prod_{j\in S}p_j=1/(|S|+1)\qquad(S\subseteq[m]).
\]
Write
\[
 u_j=1-p_j,\quad s=\sum_j u_j,\quad
 \ell=\lfloor m/4\rfloor,\quad N=m-\ell,\quad
 \alpha=\ell/m,\quad x=\alpha s,\quad\theta=1/16.
\]
For a uniform $\ell$-subset $B$, put
$A_B=\prod_{j\in B}p_j$, $\bar A=\E_B A_B$, and
$v_B=N^{-1}\sum_{j\notin B}\ell u_j$.
The probability measure $\mu$ on $x$ has row masses
$(\ell+1)w_q\bar A(q)$. The measure $\nu$ has row masses $mw_q$ and
total mass $m$. Let $\beta$ have probability density
\[
 b_\ell(z)=\frac{\ell+1}{\ell}(1-z/\ell)^\ell
                         \boldsymbol1_{[0,\ell]}(z).
\]
All constants below are absolute. We work with sufficiently large $m$;
enlarging constants covers the remaining range.

The two crucial reviewed inputs are
\begin{align}
 \max_j u_j&\le C(s+\delta)/m,
       &\delta&=m^{-1/3},\label{relative}\\
 \nu([0,t])&\le C(t+h)\quad(t\ge0),
       &h&=m^{-3/4}.\label{cdf}
\end{align}
The first is Theorem 1 of
\path{../rational_designs/asymptotic/quantitative_relative_endpoint_balance.tex}.
The second is the first bootstrap step, from exponent $1/2$ to $3/4$,
in \path{near_quadratic_per_die_bootstrap.tex}. Only this single step
is needed; no constants depending on a limiting exponent are used.

We also use the earlier exact polarization identity
\begin{equation}\label{exact}
 \beta(P)=(\ell+1)\E_w\E_B
       [A_B\Pol_N(P)((\ell u_j)_{j\notin B})]
\end{equation}
for $\deg P\le N$, and the elementary estimate
\begin{equation}\label{abar}
 \bar A\le e^{-\alpha s}=e^{-x}.
\end{equation}
For a centered real vector $y$ of length $N$, with
$V=\sum_j y_j^2$, the entire-function Cauchy estimate gives
\begin{equation}\label{centered}
 |e_k(y)|/\binom Nk\le(\sqrt{ekV}/N)^k.
\end{equation}
These facts and the Gaussian Laguerre identities are proved in
\path{../power_per_die_lower_bound.tex} and
\path{gaussian_laguerre_variance_lemmas.tex}.

\section*{An all-orders estimate for the particular endpoint tests}
Set
\[
 T_d(y)=\sum_{j<d}(1-j/d)L_j(y)=d^{-1}L_{d-1}^{[2]}(y),
\quad U_d(z)=d^{-2}T_d(\theta z)^2,
\quad R_d(z)=(2/(\theta d)-z)T_d(\theta z)^2.
\]
Square brackets denote an associated-Laguerre parameter.

\begin{lemma}\label{allorders}
For every $d\ge1$, $0\le j\le d-1$, and $z>0$,
\begin{equation}\label{Tder}
 \left|\frac{d^j}{dz^j}T_d(\theta z)\right|
 \le\frac{C e^{\theta z/2}}z
                       \left(C\sqrt{d/z}\right)^j.
\end{equation}
Higher derivatives vanish. Consequently, for $z\ge1/(2d)$ and every $k$,
\begin{align}
 \frac{|R_d^{(k)}(z)|}{k!}
 &\le\frac{C(k+1)e^{\theta z}}z
                      \frac{(C\sqrt{d/z})^k}{k!},\label{Rder}\\
 \frac{|U_d^{(k)}(z)|}{k!}
 &\le\frac{C e^{\theta z}}{d^2z^2}
                      \frac{(C\sqrt{d/z})^k}{k!}.\label{Uder}
\end{align}
All constants are uniform in the derivative order.
\end{lemma}
\begin{proof}
The off-diagonal Gaussian matrix-coefficient bound gives, for integers
$r,a\ge0$,
\[
 y^{a/2}|L_r^{[a]}(y)|
       \le e^{y/2}\sqrt{(r+a)!/r!}.
\]
Apply it with $r=d-j-1$, $a=j+2$, whose sum is always $d+1$.
Since
\[
 D_z^j T_d(\theta z)=(-\theta)^jL_{d-j-1}^{[j+2]}(\theta z)/d,
\]
and $(d+1)!/(d-j-1)!\le(d+1)^{j+2}$, it proves \eqref{Tder}.
This is valid simultaneously for every derivative order, rather than
only for fixed order. The product rule proves \eqref{Uder}.
For \eqref{Rder}, use
\[
 R_d^{(k)}=(2/(\theta d)-z)(T_d(\theta z)^2)^{(k)}
                        -k(T_d(\theta z)^2)^{(k-1)}.
\]
When $z\ge1/(2d)$, the linear factor has magnitude at most $Cz$,
and the second term costs at most an additional factor
$k/\sqrt{dz}\le\sqrt2 k$. This proves the assertion.
\end{proof}

We will sum the whole polarization expansion, using
\begin{equation}\label{series}
 \sum_{k=2}^\infty (k+1)a^k k^{k/2}/k!
                 \le C a^2 e^{Ca^2}\qquad(a\ge0).
\end{equation}
For completeness, split into even and odd orders. At $k=2j$,
$(2j)^j/(2j)!\le2^j/j!$, since
$(2j)!/j!\ge j^j$. At $k=2j+1$, the same argument bounds the
ratio by $C^{j+1}(j+1)/j!$. Summing the two exponential series gives
$C(a^2+a^3)(1+a^4)e^{Ca^2}$; increasing $C$ absorbs the extra factors
and gives \eqref{series}. This argument supplies absolute constants.

\section*{Improved subset moments and lower tails}
On rows with $s\le m/256$, the success-weighted subset law
$\Pr(B)=A_B/e_\ell(p)$ has $\Delta=v_B-x$ satisfying
\begin{equation}\label{moments}
 0\le\E\Delta\le Cx(x+\delta)/m,\qquad
 \E\Delta^2\le C\{x(x+\delta)/m+x^2(x+\delta)^2/m^2\}.
\end{equation}
Indeed the exact bias identity and negative inclusion covariances
proved in \path{../rational_designs/asymptotic/weighted_success_subset_moments.tex}
bound the bias and variance by constants times
$\sum_j u_j^2$. By \eqref{relative}, this sum is at most
$Cs(s+\delta)/m$; also $p_j\ge1/2$ here, and $x\asymp s$.

For the uniform subset law, the sharpened lower-tail bound is
\begin{equation}\label{tail}
 \Pr_B(v_B<x/2)\le\exp[-cmx/(x+\delta)].
\end{equation}
To verify it, the population $z_j=\ell u_j$ has mean $x$ and maximum
at most $C(s+\delta)$. For any upper bound $R$ on that maximum,
Maclaurin's inequality gives
\[
 \E_Be^{-Nv_B/R}\le
 \left(m^{-1}\sum_j e^{-z_j/R}\right)^N
 \le\exp[-Nx(1-e^{-1})/R].
\]
Markov's inequality gives \eqref{tail}, as in the earlier bootstrap.
When success weighting is present, the unnormalized bad contribution
is still bounded by \eqref{tail}: explicitly
$\bar A\Pr_{\rm weighted}(\mathrm{bad})
=\E_{B\,\rm uniform}A_B\boldsymbol1_{\rm bad}
\le\Pr_{\rm uniform}(\mathrm{bad})$.

\section*{Comparison at linear degree}
\begin{lemma}\label{comparison}
There is an absolute $\kappa_0>0$ such that, for each fixed
$0<\kappa\le\kappa_0$ and $d=\lfloor\kappa m\rfloor$,
\begin{align}
 |\mu(R_d)-\beta(R_d)|&\le C\kappa+o(1),\label{Rcompare}\\
 |\mu(U_d)-\beta(U_d)|&=o(1/m).\label{Ucompare}
\end{align}
The constant $C$ is independent of $\kappa$; the little-oh terms may
depend on the fixed positive $\kappa$.
\end{lemma}
\begin{proof}
Choose $\kappa_0$ small enough that $2d\le N$ and $x\le d$ implies
$s\le m/256$. We again compare exact polarization to $P(v_B)$ and
then compare $P(v_B)$ to $P(x)$.

First consider $x\le1/d$. The global Gaussian bounds, with coefficient
product $B_R=O(d)$ or $B_U=O(1)$ and variance proxy
\[
 \sum_{j\notin B}(z_j-v_B)^2\le C N(s+\delta)v_B,
\]
give full polarization error, after subset averaging, at most
$C B_P(d/m)(x+\delta)$ in this range. Indeed its series starts at
order two and \eqref{series} has argument squared
$O((d/m)(s+\delta))=o(1)$. The weighted Taylor argument using
\eqref{moments} gives the same bound. One can use the global derivative
envelopes $C B_P\sqrt{d/z}\,e^{\theta z}$ and
$C B_P(d/z)e^{\theta z}$: the factor $x$ in \eqref{moments} cancels
the singular Taylor remainder, as in the reviewed power-bound proof.
At $x=0$ all variables vanish and both errors are zero.
By \eqref{cdf}, the integrated small-$x$ contributions are therefore
\begin{align}
 E_R^{\rm small}&\le C\kappa d(1/d+\delta)(1/d+h)=o(1),\label{smallR}\\
 E_U^{\rm small}&\le C\kappa(1/d+\delta)(1/d+h)=o(1/m).\label{smallU}
\end{align}
Here $d\asymp m$, $\delta=m^{-1/3}$, $h=m^{-3/4}$, and hence
$d\delta h\to0$ and $\delta h=o(1/m)$.

Next let $1/d<x\le d$. A subset is good when $v_B\ge x/2$;
always $v_B\le x/(1-\alpha)<2x$. Combining \eqref{Rder}--\eqref{Uder}
with \eqref{centered} cancels every power $v_B^{-k/2}$.
On good subsets, \eqref{series} bounds the entire polarization error
by
\[
 C W_P(x)e^{\theta v_B}\frac{d(s+\delta)}m
                           \exp[Cd(s+\delta)/m],
 \qquad W_R(x)=1/x,\quad W_U(x)=1/(d^2x^2).
\]
Multiply by $A_B$ and average. Since $\bar A\le e^{-x}$,
$s\asymp x$, and $d/m\le\kappa$, choosing $\kappa_0$ sufficiently
small leaves a fixed exponential margin. Thus the respective
unnormalized row errors are bounded by
\begin{align}
 C\kappa e^{-cx}(1+\delta/x),\label{Renvelope}\\
 \frac{C}{dm}e^{-cx}(1/x+\delta/x^2).\label{Uenvelope}
\end{align}
These quantities are subsequently integrated against a constant
multiple of $\nu$, because $(\ell+1)w\asymp mw$.

For bad subsets use the coarser global Gaussian bounds for the whole
series. They are at most a polynomial in $m$ times $e^{Cx}$ in this
region. Multiplication by $A_B\le1$ and \eqref{tail} make the bad
contribution negligible. If $1/d<x\le\delta$, the exponent is at most
$-cm/(d\delta)$, which is $-c/(\kappa\delta)$; if $x\ge\delta$,
it is at most $-cm+Cd\le-cm/2$ after reducing $\kappa_0$.
Both are superpolynomially small. Integration against total mass $m$
preserves this conclusion.

For the weighted averaging error, keep the full linear bias
$P'(x)\E\Delta$ and split only the Taylor remainder into good and bad
subsets. On good subsets all intermediate points are in $[x/2,2x]$.
For $R_d$, the first two derivative bounds are
\[
 |R_d'|\le C\sqrt d\,x^{-3/2}e^{2\theta x},\qquad
 |R_d''|\le C d x^{-2}e^{2\theta x}.
\]
For $U_d$ they are
\[
 |U_d'|\le C d^{-3/2}x^{-5/2}e^{2\theta x},\qquad
 |U_d''|\le C d^{-1}x^{-3}e^{2\theta x}.
\]
With \eqref{moments}, and after multiplication by $\bar A$, these
give the same envelopes \eqref{Renvelope}--\eqref{Uenvelope}.
For example, the leading quadratic terms are exactly of these forms;
the ratio of the linear term to the corresponding envelope is at
most $C\sqrt{x/d}\le C$. The squared-bias term introduces only
$x(x+\delta)/m$, whose remaining polynomial factor is absorbed in
the exponential margin. On bad subsets bound the remainder by
$|P(v_B)|+|P(x)|+|P'(x)||\Delta|$. This is polynomial in $m$ times
$e^{Cx}$; the weighted normalization cancels exactly as explained
after \eqref{tail}. The same negligible bad probability applies.

For $x>d$, use the coarser global Gaussian full-series bound, with
coefficient product $B_P$. Equations \eqref{series} and \eqref{abar}
leave $e^{-cx}$ once $\kappa_0$ is small. Its integral is bounded by
a polynomial in $m$ times $e^{-cd}$. For averaging, the nonsingular
first derivative bound and uniform-subset Cauchy--Schwarz give the
same conclusion. These contributions are superpolynomially small
because $d\asymp m$.

It remains to integrate \eqref{Renvelope}--\eqref{Uenvelope}.
Stieltjes integration by parts under \eqref{cdf} gives
\begin{align*}
 \int_{1/d}^\infty e^{-cx}\,d\nu(x)&\le C,\\
 I_1:=\int_{1/d}^\infty e^{-cx}x^{-1}\,d\nu(x)
       &\le C(\log d+hd),\\
 I_2:=\int_{1/d}^\infty e^{-cx}x^{-2}\,d\nu(x)
       &\le C(d+hd^2).
\end{align*}
One may bound the cutoff endpoint term by
$\nu([0,1/d])$ times the endpoint value; the same displayed bounds
result, including possible atoms. Hence the middle $R$ error is
\[
 C\kappa(1+\delta I_1)
 \le C\kappa[1+\delta\log d+\delta hd]=C\kappa+o(1).
\]
The middle $U$ error is at most
\[
 \frac C{dm}(I_1+\delta I_2)
 \le C\left[\frac{\log d}{dm}+\frac hm+\frac\delta m
                         +\frac{\delta hd}m\right]=o(1/m).
\]
Together with \eqref{smallR}--\eqref{smallU}, the negligible tails,
and \eqref{exact}, this proves the lemma.
\end{proof}

\section*{The endpoint and its atom}
The scalar taper calculation in the reviewed power-bound proof gives
\[
 \beta(R_d)\ge1/(96\theta)>0.
\]
First choose a fixed $\kappa>0$ small enough that the $C\kappa$ in
\eqref{Rcompare} is less than half this constant, and then let $m$
be sufficiently large. Since $R_d(z)\le0$ for $z\ge2/(\theta d)$,
positivity of $\mu(R_d)$ forces some row below that threshold.
The last row minimizes $x$, so
\[
 x_K<2/(\theta d).
\]
Let $r=\lfloor d/16\rfloor$. Then $r\asymp m$ and the comparison
lemma applies with the smaller fixed proportionality constant.
For $j<r$, $j\theta x_K\le1/8$, so the elementary estimate
$|L_j(\theta x_K)-1|\le e^{j\theta x_K}-1<1/3$ gives
$U_r(x_K)\ge c>0$. On the other hand orthogonality and
$b_\ell(z)\le2e^{-\theta z}$ give $\beta(U_r)\le C/r$.
Therefore
\[
 c\mu(K)\le\mu(U_r)\le C/r+o(1/m)\le C/m.
\]
Finally $s_K=x_K/\alpha=O(1/m)$ and
$A_B(K)\ge1-s_K\ge1/2$. Thus
\[
 (\ell+1)w_K/2\le\mu(K)\le C/m,
 \qquad w_K\le C/m^2.
\]
Reversal of rows and complementation of all columns preserve the
mixed moments and give the same bound for the first row. For a
distinguished equiprobable die, every row weight is $1/K$ and
$m=n-1$, proving Theorem~\ref{main}.
\end{document}
